% Generated draft: requires mathematical and source review.
\documentclass[11pt]{article}
\usepackage[margin=25mm]{geometry}
\usepackage{fontspec}
\setmainfont{Latin Modern Roman}
\usepackage{amsmath,amssymb,mathtools,unicode-math}
\setmathfont{Latin Modern Math}
\usepackage{enumitem,needspace,xurl,hyperref}
\setlength{\parindent}{0pt}
\setlength{\parskip}{5pt}
\setlength{\emergencystretch}{3em}
\newcommand{\N}{\mathbb{N}}
\newcommand{\Z}{\mathbb{Z}}
\newcommand{\Q}{\mathbb{Q}}
\newcommand{\R}{\mathbb{R}}
\newcommand{\C}{\mathbb{C}}
\newcommand{\F}{\mathbb{F}}
\newcommand{\A}{\mathbb{A}}
\newcommand{\PP}{\mathbb{P}}
\newcommand{\E}{\mathbb{E}}
\newcommand{\eps}{\varepsilon}
\newcommand{\dd}{\,\mathrm{d}}
\begin{document}
\section*{UnsolvedMath ID: 30006297}
\textbf{Source problem code:} catalogue entry
\subsection{Definitions and mathematical statement}
Let \(X_n\) denote the ambient permutohedral variety in the supplied report, and let \(W_L\subseteq X_n\) denote the wonderful variety associated with a linear realization \(L\) of a matroid. Write \(d=\dim W_L\), let \(K_{W_L}\) denote a canonical divisor, and let \(\partial W_L\) denote the boundary divisor of the intended wonderful compactification. The supplied excerpt does not specify the ground field, the hypotheses on the realization, the precise wonderful model, or the boundary construction. Consequently, these symbols refer to the report's intended construction; the quantification over fields and realizations requires clarification before this becomes an unconditional statement.

Interpret divisors occurring inside \(\mathcal O\) as integral Cartier divisors, so that the displayed sheaves are line bundles. This interpretation should be checked against the report. A Cartier divisor \(D\) on a projective variety is nef if \(D\cdot C\geq 0\) for every integral curve \(C\). A Cartier divisor \(A\) on \(W_L\) is big if its Iitaka dimension is \(d\). The canonical-divisor notation presupposes the usual smooth wonderful model, whose construction is not supplied here.

\textbf{Nef-divisor vanishing (Conjecture 3 in the supplied excerpt).} For every nef divisor \(D\) on \(X_n\), the proposed vanishing is
\[
H^i\!\left(W_L,\mathcal O_{W_L}(D|_{W_L})\right)=0
\qquad\text{for every integer }i>0.
\]
For every such \(D\) for which \(D|_{W_L}\) is also big, the further proposed vanishing is
\[
H^i\!\left(W_L,\mathcal O_{W_L}(K_{W_L}+D|_{W_L})\right)=0
\qquad\text{for every integer }i>0.
\]
The nefness hypothesis is imposed on the ambient divisor \(D\), whereas the additional bigness hypothesis is imposed on its restriction.

\textbf{Anti-log-canonical question.} Determine the validity of
\[
H^i\!\left(W_L,\mathcal O_{W_L}(-K_{W_L}-\partial W_L)\right)=0
\qquad(0\leq i<d).
\]
These are the nonnegative cohomological degrees covered by the source's phrase \(i<\dim W_L\); when \(d=0\), there are no such degrees. The supplied cleaned statement asks this for a realizable matroid without a connectedness restriction. However, the original excerpt introduces it as a strengthening of a result assuming that the matroid of \(L\) is connected. Both formulations must be retained: the intended connectedness hypothesis for this question remains unresolved.

\textbf{Supplied partial results and qualifications.} The excerpt reports that both assertions of Conjecture 3 hold for divisors in a full-dimensional subcone \(C\) of the nef cone of \(X_n\), generated by the ``simplicial generators'' of its reference [2] (Theorem 4, attributed there to [9]). The generators and the cone are not defined in the supplied material. It describes this as a consequence of Brion's result [5] on multiplicity-free subvarieties of flag varieties, and says that [9] also contains a version for all matroids, including nonrealizable matroids, requiring additional techniques. That broader version is background and does not define \(W_L\) for a nonrealizable matroid here.

The excerpt also reports the log-canonical vanishing
\[
H^i\!\left(W_L,\mathcal O_{W_L}(K_{W_L}+\partial W_L)\right)=0
\qquad(i>0)
\]
(Theorem 5, attributed there to [8]). For a connected matroid, it reports
\[
(-1)^{r-1}\chi\!\left(W_L,\mathcal O_{W_L}(-K_{W_L}-\partial W_L)\right)\geq 0
\]
(Theorem 6, attributed there to [4]), where \(r\) is the report's rank parameter, apparently the dimension of the realization \(L\), and
\[
\chi(W_L,\mathcal F)=\sum_{j=0}^{d}(-1)^j\dim_k H^j(W_L,\mathcal F).
\]
The convention for \(r\) should be checked. The supplied excerpt explicitly describes the stronger anti-log-canonical vanishing as open. These are attributed source statements, not independently verified results or a new assessment of their present status.

The excerpt attributes the methods for Theorems 5 and 6 to tautological bundles of matroids introduced in its reference [3]. It describes a subbundle \(S_L\) of a trivial equivariant bundle, with fiber \(t^{-1}L\) over a point \(t\) of the dense projective torus, and a quotient bundle \(Q_L\). It further reports that \(W_L\) is the zero locus of a section of \(Q_L\), citing [3, Theorem 7.10], and mentions an observation implicit in [11]. The precise torus conventions and the displayed resolution cannot be recovered reliably from the supplied transcription.
\subsection{Short English statement}
For a realizable matroid's wonderful variety, the conjecture asks for higher-cohomology vanishing for restrictions of ambient nef divisors, and for their canonical twists when the restrictions are big. A further question asks whether the anti-log-canonical line bundle has no cohomology below the dimension. The ground field, compactification conventions, and connectedness scope of the last question require clarification. The supplied record labels the problem partially solved; its reported partial results and dated literature assessment have not been independently verified.
\subsection{Sources}
\begin{itemize}
\item Supplied local registry record: \texttt{data/unsolvedmath/problems.json}, permanent UnsolvedMath ID \(30006297\), source problem code \texttt{catalogue entry}, title ``Vanishing Theorems for Divisors on Matroid Wonderful Varieties.'' The record supplies the status \texttt{partially\_solved}. Its \texttt{corrected\_verified} statement label is accompanied by an explanation limiting the changes to formula delimiters, notation, grammar, and typography; it is not independent mathematical verification. No individual proposer is identified.
\item \emph{Toric Geometry} (2025), \emph{Oberwolfach Reports}, pp.\ 883--940, DOI 10.4171/owr/2025/19. Supplied source report and literature-source link: \url{https://doi.org/10.4171/owr/2025/19}. The excerpt identifies Conjecture 3 and Theorems 4--6, but does not supply their exact pages or the associated bibliography.
\item Eur, \emph{Cohomologies of tautological bundles of matroids}. Supplied literature-source link: \url{https://link.springer.com/article/10.1007/s00029-024-00979-7}. The supplied assessment describes this as a 2024 paper proving log-canonical vanishing in arbitrary characteristic. This attribution and scope have not been independently checked.
\item Literature assessment dated exactly \texttt{2026-08-22}: supplied dated curation, not independently verified. It labels the problem partially solved, reports a full-dimensional cone of nef divisors satisfying the conjecture and the log-canonical case, and states that the stronger anti-log-canonical vanishing remains open. The background expressly requests independent verification of time-sensitive status and cited claims.
\item Reference labels [2], [3], [4], [5], [8], [9], and [11] are internal to the supplied report excerpt. Only [5] is named there as a result of Brion; complete bibliographic identities are not supplied and have not been reconstructed.
\end{itemize}
\end{document}
